\section{Introduction}

Two Android packages may be related because of copying, a legitimate fork, a shared advertising library, a common template or independent generation. Similar features can support several of these histories. The practical question is therefore not only whether artifacts resemble one another, but which relationship the observed evidence can establish for a developer, marketplace or incident responder.

Android security surveys already distinguish analysis observations, detection outcomes and ecosystem roles~\cite{xuandroid,tamandroid}. Repackaging mechanisms recover code/graph reuse, designed marks and localized grafting~\cite{droidmoss,attackclones,piggybacking}; advertising and watermark evidence add other signals~\cite{adrob,appink}. TPL methods separate modules and identify components~\cite{commonlibraries,libd,libradar}, subject to shrinking, reference coverage and version ambiguity~\cite{libscout,libid,atvhunter}. Authenticated build steps and delegated release roles answer different questions about materials, actors, products and authority~\cite{intoto,tuf,diplomat}. Publication histories add another observation~\cite{chainiac}. These mechanisms are complementary; their guarantees do not compose automatically.

Generated development extends the transformation space. Text2App and AppForge concern mobile artifacts, while assistant-security, dependency and sensitive-output studies mostly concern other task populations~\cite{text2app,appforge,pearce}. Their observations motivate context and accepted-output bindings, not a present-day Android risk rate~\cite{perry,sandoval,spracklen}. Output association also differs from authenticated training or repository ancestry~\cite{codexleaks}. This review treats such workflows as an extension of Android lineage rather than a census of generative systems.

Prior synthesis provides a substantial foundation: RePack and newer repackaging reviews compare detection mechanisms~\cite{repack,wurepack2026}; TPL reviews address component stages and security~\cite{tpl,zengtpsurvey2024}; supply-chain work organizes attacks, properties and design choices~\cite{astra,properties,ladisa}. Research-direction/risk reviews and empirical SBOM studies further distinguish production records, governance and composition~\cite{williams2025,cyberrisk2026,bi2024sbom}; transit integrity and actual content are separate concerns~\cite{bomsaway2024}. AStRA's retrieved v2 places generated software in its outlook, while Williams et al. directly treat generation and repair. The comparison here is therefore about how to join existing evidence, not a claim that these subjects lacked predecessors.

We ask four research questions:
\begin{description}
\item[RQ1: Claims.] Which distinct lineage and assurance propositions do the mechanisms support?
\item[RQ2: Transformations.] How do rewriting, compilation, component reuse and generation change the validity of artifact evidence?
\item[RQ3: Composition.] Which bindings join artifact, component, process, authority and behavior observations?
\item[RQ4: Generated workflows.] Which additional causal-context and accepted-output records are needed?
\end{description}

The \CorpusWorks{}-work evidence map contains \PrimaryWorks{} primary/system studies. Its seven-field coding and unresolved records are described in Section~\ref{sec:method}; incomplete comparison cells do not establish a literature-wide gap. The synthesis supplies seven distinct claim categories, a mechanism map, counter-histories and an Android release-chain walkthrough. Artifact evidence remains useful when process records are missing; authenticated records supply recorded direction and accountability; release policies establish scoped authority; and property-specific analyses address behavior. Explicit object, predecessor, principal and policy bindings let these observations support one reviewable case without substituting one claim for another.

Sections~\ref{sec:model}--\ref{sec:ai} develop the evidence model and mechanism syntheses. Sections~\ref{sec:synthesis} and~\ref{sec:agenda} give the composition blueprint and research agenda.
